\subsection{The coherent bundles and analytic invariants of the fourfold}

We retain the complex orientation and the integral marking of P1.1.
Let $\widehat B$ be the root orbifold of orders $5,4$, let
$h:M\to\widehat B$ be the map supplied by the two multiple divisors,
and let $\pi:\widehat B\to\mathbb P^1$ be its coarse map.
Write $L=T^2$; in the pencil coordinate it is
\[
 L=\mathcal O_{\widehat B}\bigl(\tfrac45[\infty]
                  -\tfrac34[0]\bigr),\qquad \deg L=\tfrac1{20}.
\]
The cusp point has line bundle $L^{20}$, and
$T_{\widehat B}=L^9$. These identities include the stabilizer
characters: their quotients have integral divisors of degree zero
on $\mathbb P^1$. Denote the cusp by $W$, its normalization by $Y$,
and its reduced singular locus by $D$.

\begin{statement}{P1A.1}{Proved}
For every admissible member, its vertical translation bundle on
$\widehat B$ and its first coherent direct image are
\[
 E_0=F^\vee=\mathcal O\oplus L^{-1}\oplus L^{-2},\qquad
 G=R^1h_*\mathcal O_M=\mathcal O\oplus L^{-1}\oplus L^{-2}.
\]
Moreover $R^qh_*\mathcal O_M=\bigwedge^qG$ for $0\le q\le3$,
$K_M=h^*L^{-6}$, and
\[
 (h^0(\mathcal O_M),\ldots,h^4(\mathcal O_M))=(1,1,0,0,0).
\]
\end{statement}
\begin{proof}
The evaluation of the marked lattice on the lifted plane $F$
identifies the Lie bundle of the smooth tori with $F^\vee$.
P1.8 gives $F=\mathcal O\oplus L\oplus L^2$; its lift as a bundle
has the same splitting. The logarithmic toric vector fields extend
this Lie bundle over the cusp, so the identification holds on the
whole orbifold, including both finite root charts.

On the smooth locus the constant-cocycle description of torus
cohomology gives $G=V/F$. In the rank-five orthogonal quotient,
$K=\ker(\alpha:P\to\mathcal O)=L\oplus L^2$ is the radical of
the restricted pairing on $P$. Indeed every wedge containing the
fixed line pairs to zero with another such wedge except through
their contraction components; removing that scalar component
leaves exactly $K$. The nondegenerate pairing consequently
identifies $W/P$ with $K^\vee$. Thus
\[
 0\longrightarrow\mathcal O e\longrightarrow V/F
   \longrightarrow L^{-1}\oplus L^{-2}\longrightarrow0.
\]
For every integer $j$, descent by invariant local Laurent series gives
\begin{equation}\label{eq:P1Adegree}
 \pi_*L^j=\mathcal O_{\mathbb P^1}(d(j)),\qquad
 d(j)=\lfloor4j/5\rfloor+\lfloor-3j/4\rfloor.
\end{equation}
The extension group $H^1(L\oplus L^2)$ therefore vanishes, so this
sequence splits as a holomorphic vector-bundle sequence.

We verify that this quotient is the coherent extension at the cusp.
The normalization resolution of $\mathcal O_W$ has one toric
threefold, seven toric surfaces, twelve toric curves and six points.
All positive coherent cohomology of those smooth complete toric
strata vanishes. Their constants form precisely the cellular
cochain complex of the dual three-torus, with the incidences in
\path{code/part1/cusp_complex.json}. Its cohomology is
$\bigwedge^q\mathbb C^3$, and its cup product is the exterior product.
The invariant flat cohomology $\ker N$ maps to these classes by
constant Cech cocycles on the toroidal covering. A component
translation in that cover evaluates the corresponding lattice
coordinate. In the nilpotent logarithmic frame, $\ker N$ maps
isomorphically to $V/F$, since $N:F\to\ker N$ is an isomorphism.
Thus these are the actual extending quotient classes, rather than
an unrelated basis of the same dimension.

Flatness over each smooth root disc follows because its base is a
DVR and the smooth total space is torsion free over it. Properness
was proved in P1.9--P1.10. The coherent fibre dimensions are the
constant binomial dimensions on smooth tori, finite root fibres,
and the cusp. Grauert's \S7 Satz~5, with these checked hypotheses,
gives local freeness and base change. The cup-product maps are
isomorphisms on every fibre, including the preceding central
cellular model, so they identify the remaining direct images with
$\bigwedge^qG$.

There is an exact sequence of locally free logarithmic bundles
\begin{equation}\label{eq:P1Alogcotangent}
0\longrightarrow h^*L^{11}\longrightarrow\Omega_M^1(\log W)
 \longrightarrow h^*F\longrightarrow0.
\end{equation}
In a chart $q=x_0\cdots x_r$, its first arrow sends
$dq/q$ to $\sum dx_i/x_i$ and its quotient is the height-zero
logarithmic cotangent bundle of the torus. These formulas also
prove exactness on every stratum. Taking determinants and
cancelling $\mathcal O_M(W)=h^*L^{20}$ yields
$K_M=h^*L^{11+3-20}=h^*L^{-6}$.

For any bundle which is a sum of powers of $L$, projection and
the exact finite averaging projector compute its direct images
by tensoring with $\bigwedge^qG$ and applying
\eqref{eq:P1Adegree}. All exponents $-1,\ldots,-5$ have coarse
degree $-1$, while exponent zero has degree zero. The Leray
sequence on the curve then gives the asserted structure-sheaf
cohomology. This includes the nonreduced coarse fibres, since
the computation was made before exact finite descent.
\end{proof}

\begin{statement}{P1A.2}{Proved}
The complete Hodge table of every admissible fourfold is the
following, with rows indexed by $p$ and columns by $q$:
\[
\begin{array}{c|ccccc}
h^{p,q}&0&1&2&3&4\\\hline
0&1&1&0&0&0\\
1&0&2&1&0&0\\
2&0&1&6&1&0\\
3&0&0&1&2&0\\
4&0&0&0&1&1
\end{array}
\]
\end{statement}
\begin{proof}
All maps needed for the two nontrivial rows can be evaluated at
the cusp. We give their matrices and the reason evaluation
determines their global ranks. In the algebraic splitting of
P1.15 write the three generators of $F$ as $(f_0,f_1,f_2)$ of
weights $(0,1,2)$, and those of $G$ as $(g_0,g_1,g_2)$ of
weights $(0,-1,-2)$. The residue of the quotient connection is
\[
R:F_0\longrightarrow G_0,\qquad
R=\begin{pmatrix}0&A&B\\C&0&D_0\\E&D_0&0\end{pmatrix},
\]
where
\begin{align*}
A&=-125(192z-3125)/12288,&
B&=-78125(96z+3125)/393216,\\
C&=-125/512,&E&=78125/32768,\\
D_0&=-146484375/2097152.
\end{align*}
Here the lifted first vector is $w_0+(4/3)ze$.
These entries result from the exact Gauss--Manin reduction of
P1.15, with the quotient projection taken after lifting $F$.
They are independently recomputed in
\path{code/part1/analytic_cohomology.py}.

The connecting map of \eqref{eq:P1Alogcotangent}, on relative
cohomology, sends $f\otimes\eta$ to $R(f)\wedge\eta$ at the
cusp. To see the actual extension class, trivialize logarithmic
forms on the toric covering. A deck translation changes a
relative logarithmic form by its period derivative times
$dq/q$. The constant Cech class of this derivative is exactly
the residue of the quotient connection. This proves the
formula with the indicated frames. It also proves its exterior
power version, with the usual alternating signs.

The floor formula gives the following absolute cohomology:
\[
\begin{array}{c|ccccc}
&h^0&h^1&h^2&h^3&h^4\\\hline
h^*F&1&3&2&0&0\\
h^*L^{11}&0&2&3&1&0\\
h^*(\bigwedge^2F)&0&2&3&1&0\\
h^*(L^{11}\otimes F)&2&6&7&3&0.
\end{array}
\]
Every group in this display comes from global sections of the
indicated direct image; none comes from $H^1$ of a line on the
coarse base. Every contributing summand has coarse degree zero,
so evaluation at the cusp is injective on its space of sections.
For the equal-dimensional source and target spaces used below
it is an isomorphism onto the displayed weight subspaces.
Thus fibrewise evaluation computes the actual ranks of the
absolute connecting maps.

For one-forms, the three middle source classes are $f_i\otimes g_i$.
In the ordered target $(g_0g_1,g_0g_2,g_1g_2)$, the matrix is
\[
C_1=\begin{pmatrix}-C&A&0\\-E&0&B\\0&-D_0&D_0\end{pmatrix},
\qquad
\det C_1=
\frac{476837158203125}{281474976710656}(384z+3125).
\]
The degree-zero connecting map is nonzero because $R(f_0)$
has nonzero $g_1,g_2$ coefficients. The degree-two connecting
map, from the classes $f_1g_0g_1,f_2g_0g_2$ to $g_0g_1g_2$,
has rank one because $D_0\ne0$.

We must exclude the sole possible zero of $\det C_1$ on the
whole admissible parameter domain. P1.15 identifies
$z=-(625\pi/128)c$ on the real slice in its algebraic splitting.
The exact exceptional parameter is therefore $c=5/(3\pi)<1$.
The certified continuation in
\path{code/part1/certified_lift_bound.py} proves
\[
h_b(1)\in[-1.03427662241632179949424884537834605259496477790218821132
          \ \mathbin{\pm}\ 8.11\cdot10^{-57}],
\]
so $h_b(1)+5/(3\pi)<0$. Since $h_b\to+\infty$ at the cusp,
the intermediate value theorem makes this parameter
nontransverse somewhere on the base. For the full complex
parameter the determinant has the same single zero, and the
same argument excludes it. The certification propagates complex
balls and adds an explicit Cauchy bound for every omitted
Taylor tail; no midpoint replacement discards an error radius.
Its exact connection and all step bounds are retained with the
program. Consequently the three connecting ranks are $1,3,1$.
The long exact sequence of \eqref{eq:P1Alogcotangent} gives
\[
h^q(\Omega_M^1(\log W))=(0,1,1,0,0).
\]

The residue sequence
\[
0\to\Omega_M^1\to\Omega_M^1(\log W)\to\nu_*\mathcal O_Y\to0
\]
is exact even at self-intersections: in a local branch chart it
records independently the coefficient of each $dx_i/x_i$ on
$x_i=0$. The normalization $Y$ is smooth complete toric, so its
only coherent structure cohomology is $H^0=\mathbb C$.
The one-form row follows immediately.

For two-forms take the second exterior power of
\eqref{eq:P1Alogcotangent}. Its short exact sequence has left
term $h^*(L^{11}\otimes F)$ and right term
$h^*(\bigwedge^2F)$. In degree two the three source classes
are $f_if_j\otimes g_ig_j$, $i<j$. The target is
$F\otimes g_0g_1g_2$. Its connecting matrix is
\[
C_2=\begin{pmatrix}-D_0&D_0&0\\E&0&-B\\0&-C&A\end{pmatrix},
\qquad \det C_2=-\det C_1.
\]
Thus it is an isomorphism, and
$h^3(\Omega_M^2(\log W))=1$.

For completeness the ordinary/logarithmic comparison in this
degree is the exact residue resolution
\[
0\to\Omega_M^2\to\Omega_M^2(\log W)
 \to\nu_*\Omega_Y^1(\log\partial Y)
 \to\nu^{[2]}_*\mathcal O_{W^{[2]}}\to0.
\]
The last space is the disjoint normalization of the seven
double surfaces. Local exactness is verified by writing a
logarithmic two-form as a sum of
$a_{ij}\,dx_i/x_i\wedge dx_j/x_j$ and terms with ordinary
differentials. The two residues of an $ij$ term have opposite
sign; their cancellation is the sole compatibility. Every local
double-surface coefficient lifts to one branch, which proves
surjectivity at the last term.
Both final sheaves have zero positive cohomology: the first
is the pushforward of the trivial rank-three toric logarithmic
bundle, and the second has complete toric surface components.
Writing $Q=\Omega_M^2(\log W)/\Omega_M^2$, we obtain
$H^q(Q)=0$ for $q\ge2$. Hence
$h^3(\Omega_M^2)=1$.
In degree zero the last residue map is zero: the fourteen toric
boundary divisors of $Y$ occur in seven opposite pairs; their
primitive normal vectors have opposite signs and their gluing
is a torus translation. This gives the zero $7$-by-$3$ residue
matrix on global logarithmic one-forms. Thus $H^0(Q)=\mathbb C^3$.
The two global logarithmic two-forms come from the
$L^{12},L^{13}$ summands of $L^{11}\otimes F$; their residues
are the independent $f_1,f_2$ directions at the cusp.
Therefore $h^0(\Omega_M^2)=0$.

Serre duality on this compact smooth complex fourfold gives
$h^{2,1}=h^{2,3}=1$. Finally the analytic Riemann--Roch theorem
in Hirzebruch, Theorem~25.3.1 and its explicit derivation
in \S25.4(a), pp.~188--190, applies to the Dolbeault elliptic
operator of $\Omega_M^p$. No Kähler hypothesis occurs there.
Using $c_1=c_3=0$, the Chern-root product gives
\[
\chi_y(M)=
\frac{(1+y)^4}{240}\,c_2^2+
\left(\frac{(-1+10y-y^2)^2}{144}
       -\frac{(1+y)^4}{120}\right)c_4
=-y+4y^2-y^3.
\]
The program verifies this polynomial exactly. It gives
$\chi(\Omega_M^2)=4$, hence $h^{2,2}=6$.
P1A.1 and Serre duality supply all other entries of the table.
\end{proof}

\begin{statement}{P1A.3}{Proved}
For every admissible fourfold,
\[
 (h^0(T_M),h^1(T_M),h^2(T_M),h^3(T_M),h^4(T_M))
 =(1,1,1,2,0).
\]
The complex parameter family has a nonzero Kodaira--Spencer map
at every member and supplies a smooth one-dimensional miniversal
Kuranishi germ.
\end{statement}
\begin{proof}
The critical ideal of $h$ is $\mathcal I_D$. Locally
$q=x_0\cdots x_r$, its generators are the products omitting
one factor. A vector field in the kernel of $dq$ has coefficients
$a_i=x_i b_i$ in the crossing directions and
$\sum b_i=0$; the remaining ordinary directions are free.
Thus the kernel is the translation bundle $h^*E_0$, and
\[
0\to h^*E_0\to T_M\to J\to0,
\qquad J=\mathcal I_D\otimes h^*L^9.
\]
The conductor square gives
\[
0\to\mathcal O_W\to\nu_*\mathcal O_Y\oplus\mathcal O_D
 \to\nu_*\mathcal O_{\partial Y}\to0.
\]
The boundary is connected and anticanonical in the toric
threefold $Y$. Its coherent cohomology is $(1,0,1)$, by
$0\to K_Y\to\mathcal O_Y\to\mathcal O_{\partial Y}\to0$
and Serre duality. The singular locus $D$ is connected and
has dimension two, as is immediate from the connected
double-face incidence graph. The map on constants is
$(a,b)\mapsto a-b$ and is onto. In higher degree the conductor
sequence makes $H^1(\mathcal O_W)\to H^1(\mathcal O_D)$
an isomorphism. The map
$H^2(\mathcal O_{\partial Y})\to H^3(\mathcal O_W)$ is onto,
since $H^3(\mathcal O_D)=H^3(\mathcal O_Y)=0$; both sides
have dimension one, so it is an isomorphism. Consequently
$H^2(\mathcal O_W)\to H^2(\mathcal O_D)$ is also an
isomorphism and $h^q(\mathcal O_D)=(1,3,3)$.

Projection and \eqref{eq:P1Adegree} give
\[
h^q(h^*E_0)=(1,1,0,0,0),\qquad
h^q(h^*L^9)=(1,2,1,0,0).
\]
The maps from the latter groups to $H^q(\mathcal O_D)$
are evaluation at the cusp followed by the just-proved
restriction isomorphisms. Each contributing coarse summand
has degree zero, so these evaluations are injective, with
ranks $1,2,1$. The critical-ideal sequence gives
$h^q(J)=(0,0,1,2,0)$. The tangent dimensions and the
isomorphism $H^1(h^*E_0)\to H^1(T_M)$ follow.

Differentiating the actual periods gives, on the universal
torus covering, the Cech cocycle of constant translations
obtained by differentiating each lattice period. Under the
constant-cocycle description of fibre cohomology this is
the map $F\to V/F$ sending $f$ to $\alpha(f)e$.
It is nonzero and occupies precisely the trivial summand
of $E_0\otimes G$. The Leray edge map identifies this summand
with $H^1(h^*E_0)$, so the global Kodaira--Spencer class is
nonzero. The finite and toroidal charts depend holomorphically
on the parameter and supply the actual compact family,
so this cocycle is the derivative of that family on every chart.

Douady's analytic versality theorem, VII.8, applies to the
compact complex manifold. The Kuranishi germ embeds in its
one-dimensional tangent space. A classifying map from the
actual parameter disc has nonzero derivative, hence its
composite with that embedding is an invertible one-variable
analytic germ. Its pullback on $\mathbb C\{t\}$ is injective,
so the defining ideal of the Kuranishi germ is zero, including
nilpotents. This proves smoothness and miniversality.
\end{proof}

\begin{statement}{P1A.4}{Proved}
The fourfold has algebraic dimension one and
\[
\Aut^0(M)=\mathbb C,\qquad
\Aut(M)=\mathbb C\times\mathbb Z/2.
\]
In a fixed holomorphic splitting, two members of the full
admissible complex parameter half-plane are biholomorphic
exactly when their parameters are equal.
\end{statement}
\begin{proof}
An effective divisor dominating the base would restrict to
a nonempty effective divisor on a general complex torus.
Pairing its Chern class with the square of a Kähler form
on that torus is strictly positive. This contradicts
$H^2(M;\mathbb Z)=0$. Thus every divisor is vertical.
The pole divisor of a meromorphic function is vertical;
its restriction to a general compact fibre is holomorphic
and therefore constant. The graph and properness of $f$
descend it to a meromorphic function on $\mathbb P^1$.
Conversely pullbacks from the base occur. Hence the
meromorphic function field is $\mathbb C(t)$ and $a(M)=1$.
Every automorphism consequently induces a Möbius map on
the base. The two distinct multiplicities and the cusp
identify all three values separately, so this map is the
identity.

On smooth fibres an automorphism is affine. Its integral
linear part is locally constant and commutes with $M_1,M_2$.
Solving the two exact matrix equations in
\path{code/part1/analytic_automorphisms.py} gives the rational
centralizer $\mathbb Q I$; its integral invertible elements
are $\pm I$. At the order-five filling an affine lift with
linear part $\epsilon I$ would have to conjugate the same
finite affine generator. For its translation $b$ this
requires
\[
(I-M_1)b=(1-\epsilon)v_1/5\pmod\Lambda.
\]
Applying the invariant integral covector $\ell$ makes the
left side zero. For $\epsilon=-1$ the right side is
$2/5\pmod\mathbb Z$, which is nonzero. Thus every
automorphism is a translation.

Let $\mathcal A$ be the sheaf of translations extending to
the specified cusp model and finite root charts. There is
an exact analytic sheaf sequence
\[
0\to j_*\Lambda\to E_0\to\mathcal A\to0
\]
on the root orbifold. Here $j$ omits the cusp and the stalk
of $j_*\Lambda$ there is $\ker N$.
To check the last surjectivity at the cusp, lift a translation
to the toroidal cover. It preserves the fan precisely when
its valuation is an integral translation of the periodic
fan. Subtract the corresponding period; the remaining
characters are nonvanishing holomorphic functions on a
disc and possess logarithms. These give a local section
of $E_0$. The kernel consists exactly of the invariant
integral periods. This also proves that the sequence
accounts for all extendable translations, including
possible component permutations.

The cohomology of $j_*\Lambda$ in degree one is calculated
from three positive meridians. For the finite points use
the integer kernels of $I+M_i+\cdots+M_i^{m_i-1}$;
for the cusp use the primitive image of $N$.
Writing these column bases as $L_1,L_2,L_0$, the cocycle
and coboundary matrices are
\[
D_1=(L_1\ M_1L_2\ M_1M_2L_0),\qquad
D_0=\begin{pmatrix}L_1^{-1}(M_1-I)\\
L_2^{-1}(M_2-I)\\L_0^{-1}(M_0-I)\end{pmatrix}.
\]
Each inverse here denotes the unique coordinates in the
displayed integral column basis, not a square matrix
inverse. The exact $6$-by-$11$ and $11$-by-$6$ integer
matrices, their bases and Smith operations are exported
in \path{code/part1/analytic_automorphisms.json} and the
underlying parabolic program. They satisfy $D_1D_0=0$ and
give $\ker D_1/\im D_0=\mathbb Z/2$. The cocycle relations
are the finite-order relations and the product of the
three meridians; the cusp condition is precisely the
vanishing of its restriction modulo $(M_0-I)\Lambda$.
Thus this complex computes the required orbifold sheaf
cohomology, with its finite stabilizer torsion retained.

Now $H^0(j_*\Lambda)=0$ by P1.1,
$H^0(E_0)=\mathbb C$, and $H^1(E_0)=0$ by the floor
calculation. The exact sequence gives
\[
0\to\mathbb C\to H^0(\mathcal A)\to\mathbb Z/2\to0.
\]
The half-period
$\tau=(0,1,1,1,1,1)^t/2$ is fixed modulo $\Lambda$
by every monodromy matrix. It extends over the finite
quotients because the translations commute there.
At the cusp $N\tau$ is integral; primitiveness of $N$
provides $\lambda\in\Lambda$ with $N\lambda=N\tau$.
The representative $\tau-\lambda\in\ker N$ is a bounded
angular half-period and extends on every toric chart.
It is a nonzero order-two global section. Since the
additive group $\mathbb C$ has no nonzero torsion, it
splits the displayed exact sequence. Translations
commute, giving the asserted full group.

For an isomorphism between two members, the same intrinsic
base argument applies. Its integral linear part still
lies in the common centralizer. Either scalar $\pm I$
preserves a complex period plane exactly by equality
of the two planes. The graphs
$(\sigma+\xi\alpha e)(P)$ are equal precisely when their
parameters $\xi$ agree, since $e\notin F$ and $\alpha$
is nonzero. Translation does not change this condition.
This proves the parameter assertion.
\end{proof}
